\section*{Declarations}
\label{sec:declarations}

\subsection*{Ethics and containment}\label{ethics-and-containment}

Every sample analysed in this work is live Windows PE malware. The \fact{baseline-n}-sample cohort and the \fact{drift-n}
dated samples of the withdrawn drift study were obtained from MalwareBazaar (abuse.ch) under its
terms of use, which permit research redistribution of samples but not of the corpus as a package.

Detonation happens only inside a CAPE sandbox on a dedicated Windows guest image reverted to a clean
snapshot between analyses. The sandbox host is reachable from one network segment and from no public
route, holds no credentials for any other system, and shares no filesystem with the host that runs
the pipeline. Binaries are held only in a directory excluded from the host's real-time scanner, are
never committed to version control, and are not redistributed with this paper. We note a limitation
of that arrangement rather than claiming it sufficient: the guest reaches the public internet during
detonation, because a sample that cannot resolve its command-and-control infrastructure behaves
differently from one that can and the network evidence channel would otherwise be empty for every
sample. Live detonation with egress is the standard arrangement for behavioural sandboxing, and it
is the arrangement whose consequences we report, including that the domains every sample contacted
turn out to be the guest image describing itself.

No user study was conducted, no personal data was collected, and no human analyst scored any output.
That last point is a limitation of the evaluation, stated as such in Threats to Validity, rather
than a claim about ethics.

\subsection*{Data availability}\label{data-availability}

The evaluation harnesses, the derivation of every number in this paper, the per-sample records they
read and the scripts that build the figures and assemble the document are released with the paper.
The archive holds every per-sample record retained by every study reported here, in the JSON the
harness wrote, together with the offline re-analysis that recomputes every interval, design effect,
minimum detectable effect and adjusted p-value from those records and its random seed, the ATT\&CK
ground-truth fixtures derived from the public \texttt{mitre-attack/attack-stix-data} repository
under the ATT\&CK Terms of Use, and the model digest and engine commit needed to reproduce the local
arm.

Four things are withheld. Malware binaries are not released and are obtainable from the sources
named in Background by the SHA-256 digests in the archive. Sandbox reports are not released in full
because they embed strings and memory contents from the samples. The sandbox host's address appears
nowhere in the archive and is referred to throughout as a private address. And one dataset used for
retrieval-corpus construction is redistributed by its authors under terms that do not permit our
redistribution, so it is referenced rather than included.

Not available. The \fact{drift-n}-sample temporal-drift study's per-sample outputs do not exist. They were
written to a path that was valid on the machine the harness was first written on and silently
relative on the machine it ran on; the study is withdrawn on that account and the withdrawal is
reported in Results rather than quietly repaired. The manifest of the \fact{drift-n} samples survives and is
included, so the cohort can be reconstructed even though the results cannot.

\subsection*{Reproducibility}\label{reproducibility}

Two runs of the offline analysis over the released records produce byte-identical output, asserted
by a test in the released suite rather than claimed here, and every interval carries the seed, the
iteration count and the number of clusters it was computed from in the artefact itself, so an
interval cannot be read without what is needed to reproduce it. The live arms are not
bit-reproducible and we do not claim they are. Decoding is at temperature \fact{decode-temperature}
on the local arm and at each provider's default on the hosted ones, one of which accepts a
reasoning-configuration parameter and does not act on it. Anyone re-running them should expect the
numbers to move and should read the intervals accordingly.

\subsection*{Use of generative AI}\label{use-of-generative-ai}

Large language models are the object of study in this work, so the pipeline under evaluation is
built from them and every measurement reported here is a measurement of their behaviour in it.
Separately, generative AI was used as an authoring and engineering aid, for drafting and revising
prose, writing evaluation harnesses and their tests, and performing literature searches whose
results were then verified by fetching each source. The verification is not incidental to this
disclosure: of the nine citations added by those searches and checked one by one, three said
something the source did not, one quoted for a sentence absent from the paper, one cited for a
practice it does not describe, and one blending two incompatible published accounts into a single
attribution. All three are corrected in the text and the corrections are recorded, because a search
summary that reads as a citation is the same class of failure this paper is about. An internal
readability assessment of report narratives was performed with a language model as a development
aid; its results are deliberately excluded and no LLM scored any result reported here. The authors
take responsibility for all content, including every number, every citation, and every claim about
what the measurements support.

\subsection*{Competing interests}\label{competing-interests}

The authors declare no competing financial or non-financial interests. No external funding
supported this work. All measurements were taken on hardware owned by the authors, except for calls
to three commercially hosted inference endpoints, which were paid for at their public list prices
and whose providers had no involvement in, or advance knowledge of, this evaluation.
